% FST-DE_DarkEnergy_Skeleton_v1_en.tex
% Skeleton draft -- Dark Energy as Residual Vacuum Free Energy
% Status: SKETCH (v1)
% Date: 2026-03-10

\documentclass[12pt,a4paper]{article}
\usepackage[utf8]{inputenc}
\usepackage[T1]{fontenc}
\usepackage{lmodern}
\input{glyphtounicode}
\pdfgentounicode=1
\usepackage{amsmath,amssymb,amsthm}
\usepackage{array}
\usepackage{tabularx}
\usepackage{booktabs}
\usepackage{xcolor}
\usepackage{graphicx}
\usepackage{geometry}
\geometry{a4paper, margin=2.5cm}
\usepackage[english]{babel}
\usepackage{xurl}
\usepackage{microtype}
\def\HyperDestNameFilter#1{en.#1}
\usepackage[unicode]{hyperref}
\hypersetup{
  colorlinks=true,
  linkcolor=blue!70!black,
  citecolor=green!50!black,
  urlcolor=blue!70!black,
  pdftitle={Dark Energy as Residual Vacuum Free Energy: A Thermodynamic Bound on the Cosmological Constant},
  pdfauthor={Lukas Geiger},
  pdfsubject={Dark energy and cosmological constant preprint},
  pdfkeywords={dark energy, cosmological constant, vacuum energy, scalar-tensor gravity, f(R)}
}
\setlength{\emergencystretch}{2em}
\newcolumntype{Y}{>{\raggedright\arraybackslash}X}

\theoremstyle{definition}
\newtheorem{definition}{Definition}[section]
\theoremstyle{plain}
\newtheorem{theorem}[definition]{Theorem}
\newtheorem{lemma}[definition]{Lemma}
\newtheorem{proposition}[definition]{Proposition}
\newtheorem{corollary}[definition]{Corollary}
\newtheorem{axiom}[definition]{Axiom}
\theoremstyle{remark}
\newtheorem{remark}[definition]{Remark}

% --- Custom commands ---
\newcommand{\Leff}{\Lambda_{\mathrm{eff}}}
\newcommand{\LUV}{\Lambda_{\mathrm{UV}}}
\newcommand{\LIR}{\Lambda_{\mathrm{IR}}}
\newcommand{\rhov}{\rho_{\mathrm{vac}}}
\newcommand{\dd}{\mathrm{d}}

\title{\textbf{Dark Energy as Residual Vacuum Free Energy:\\
A Thermodynamic Bound on the Cosmological Constant}\\
\large\textit{Part of the FST Programme~\cite{GeigerFST-Hub-2026}}}

\author{Lukas Geiger\\
\small Independent Researcher, Bernau, Germany\\
\small ORCID: \href{https://orcid.org/0009-0005-7296-1534}{0009-0005-7296-1534}}

\date{July 2026}

\begin{document}
\maketitle
\vspace{-3.2em}

% ============================================================
\begin{abstract}\small
This revision audits the mathematical core of the proposed thermodynamic
dark-energy framework.  For the dimensionally closed reduced function
\(\Phi_{\rm red}=A+B\rho\ln(\rho/\rho_{\rm Pl})+C\rho^2\), \(B,C>0\),
existence, uniqueness, strict convexity, and an explicit Lambert-\(W\)
minimiser are proven.  The original momentum-space expression is
dimensionally inhomogeneous and does not determine the matching coefficients.
Consequently the scaling
\(\Lambda_{\rm eff}\sim H_0^2/[c^2\ln(M_{\rm Pl}/H_0)]\) remains an external
RG-matching ansatz.  Its corrected numerical value is
\(3.8\times10^{-55}\,{\rm m}^{-2}\), about \(290\) below observation.

The scalar-sector audit also changes the physical status.  The
P\"oschl--Teller potential \(V_0\,{\rm sech}^2(\phi/\phi_0)\) has an unstable
maximum at \(\phi=0\); the earlier tanh history is a phenomenological proxy,
not a derived stable attractor.  Standard quasistatic metric \(f(R)\)
perturbations approach a slip of \(1/2\) (or reciprocal \(2\), by convention)
and a GR-like lensing combination \(\Sigma\simeq1/F\), not \(5/4\) and \(9/8\).
Finally, the Hu--Sawicki environment ledger yields correct necessary
curvature thresholds but no source-bound local profile, so neither
\(|f_{R0}|=1.2\times10^{-5}\) nor \(10^{-6}\) is by itself a Cassini proof.
The framework is therefore a conditional research programme, not a closed
solution of the cosmological-constant or screening problems.
\vspace{1.5em}
\noindent\textbf{Keywords:} cosmological constant, dark energy, vacuum energy,
scalar-tensor gravity, gravitational slip
\end{abstract}

\clearpage
% ============================================================
\section{Introduction}\label{sec:intro}

This revision records the mathematical status after a line-by-line proof/paper
audit on 16 July 2026.  The purpose is narrower than that of the earlier
skeleton: it separates the statements that follow from a dimensionally closed
one-variable model from the physical identifications that remain open.

The audit changes four headline claims.  First, the original momentum-space
integrand was dimensionally inhomogeneous and therefore did not define a
physical free-energy density without explicit matching coefficients.  Second,
the proposed P\"oschl--Teller saturation law was not derived from the canonical
Klein--Gordon--Friedmann system and its hilltop is unstable.  Third, the
thermodynamic beta function defined as the derivative of a stationarity
condition is identically zero.  Fourth, standard metric \(f(R)\) perturbation
theory does not predict the previously stated values \(\eta=5/4\) and
\(\Sigma=9/8\).  These claims are withdrawn below.

\section{Dimensionally closed variational model}\label{sec:functional}

\subsection{Raw motivation and its limitation}

The earlier motivating expression was
\begin{equation}\label{eq:Phi-raw}
 \Phi_{\rm raw}(\rho)=
 \int_{\LIR}^{\LUV}
 \left[
 \frac{\omega(k)}{2}
 +T_{\rm dS}\rho\ln\!\left(\frac{\rho}{\rho_{\rm Pl}}\right)
 +\frac{8\pi G\rho^2}{k^2}
 \right]\frac{4\pi k^2}{(2\pi)^3}\,\dd k .
\end{equation}
In natural units, the three terms inside the square brackets have dimensions
\(E\), \(E^5\), and \(E^4\), respectively.  They cannot be added as written.
Planck-unit rescaling does not repair the problem: it hides powers of
\(M_{\rm Pl}\) and \(H\), precisely the powers that determine the proposed
cosmological scaling.  Equation~\eqref{eq:Phi-raw} is therefore only a
schematic motivation, not a physical functional and not a derivation from a
mode sum.

\subsection{Reduced model}

\begin{definition}[Reduced free-energy density]\label{def:functional}
For \(\rho>0\), let
\begin{equation}\label{eq:Phi}
 \Phi_{\rm red}(\rho)
 =A+B\,\rho\ln\!\left(\frac{\rho}{\rho_{\rm Pl}}\right)+C\rho^2 ,
 \qquad B>0,\quad C>0 .
\end{equation}
Here \([\rho]=[\Phi_{\rm red}]=E^4\), \([A]=E^4\), \([B]=1\), and
\([C]=E^{-4}\).  The constants \(B\) and \(C\) are renormalised matching
coefficients.  Their values are not supplied by Eq.~\eqref{eq:Phi-raw}.
\end{definition}

The kernel
\begin{equation}\label{eq:Vgrav}
 V_{\rm grav}(\rho,k)=\frac{8\pi G\rho^2}{k^2}
\end{equation}
has the dimension of an energy density and can motivate a positive quadratic
term after an explicitly specified integration and matching prescription.
Dimensional analysis selects it as a lowest-order monomial only after the
assumptions of one power of \(G\), quadratic dependence on \(\rho\), and
inverse powers of \(k\) are imposed.  It does not establish uniqueness among
general kernels, and a quasistatic \(4/3\) force enhancement does not by itself
derive an off-shell free-energy functional.

\section{Proven result and conditional scaling}\label{sec:main}

\begin{theorem}[Unique minimiser of the reduced model]\label{thm:main}
The function \(\Phi_{\rm red}\) in Definition~\ref{def:functional} has exactly
one global minimiser \(\rho_\ast>0\).  It satisfies
\begin{equation}\label{eq:stationarity}
 B\left[1+\ln\!\left(\frac{\rho_\ast}{\rho_{\rm Pl}}\right)\right]
 +2C\rho_\ast=0
\end{equation}
and is given explicitly by
\begin{equation}\label{eq:rho-lambert}
 \rho_\ast=\frac{B}{2C}\,
 W\!\left(\frac{2C\rho_{\rm Pl}}{eB}\right),
\end{equation}
where \(W\) is the principal Lambert function.
\end{theorem}

\begin{proof}
The first derivative tends to \(-\infty\) as \(\rho\to0^+\) and to
\(+\infty\) as \(\rho\to\infty\).  Moreover,
\begin{equation}
 \Phi_{\rm red}''(\rho)=\frac{B}{\rho}+2C>0 .
\end{equation}
Thus the derivative has exactly one zero and strict convexity makes that zero
the unique global minimiser.  Rearranging Eq.~\eqref{eq:stationarity} gives
\((2C\rho_\ast/B)\exp(2C\rho_\ast/B)=
2C\rho_{\rm Pl}/(eB)\), which yields Eq.~\eqref{eq:rho-lambert}.
\end{proof}

The associated curvature convention is
\begin{equation}\label{eq:Leff}
 \Leff=\frac{8\pi G}{c^4}\rho_\ast .
\end{equation}
No dependence on \(H_0\) follows until \(B/C\) is independently matched.
The frequently used expression
\begin{equation}\label{eq:bound}
 \Leff^{\rm ansatz}
 \sim\frac{H_0^2}{c^2\ln(M_{\rm Pl}/H_0)}
\end{equation}
is therefore a conditional RG-matching ansatz, not a consequence of
Theorem~\ref{thm:main}.

\begin{corollary}[Arithmetic of the conditional ansatz]\label{cor:numerical}
For \(H_0=67.4\,{\rm km\,s^{-1}\,Mpc^{-1}}\) and
\(L=\ln(M_{\rm Pl}/H_0)\simeq140\),
\begin{equation}
 \frac{H_0^2}{c^2L}\simeq3.8\times10^{-55}\ {\rm m}^{-2}.
\end{equation}
Compared with
\(\Lambda_{\rm obs}\simeq1.11\times10^{-52}\ {\rm m}^{-2}\)
from Planck~2018~\cite{Planck2018}, the ansatz is low by a factor of about
\(2.9\times10^2\).  This is an unresolved normalisation/matching factor, not
order-of-magnitude agreement.
\end{corollary}

\subsection{Scalar-sector status}\label{sec:scalar-tensor}

Consider the phenomenological Lagrangian
\begin{equation}\label{eq:Lagrangian}
 \mathcal L=
 \frac{R+\gamma R^2}{16\pi G}
 +\frac12(\partial\phi)^2
 -V_{\rm PT}(\phi)+\mathcal L_m ,
 \qquad
 V_{\rm PT}(\phi)=V_0\,{\rm sech}^2(\phi/\phi_0).
\end{equation}
For the quadratic curvature sector the scalaron mass is
\(m_s^2=1/(6\gamma)\).  The scalar potential obeys
\begin{equation}\label{eq:VPT-curvature}
 V_{\rm PT}'(0)=0,\qquad
 V_{\rm PT}''(0)=-\frac{2V_0}{\phi_0^2}<0 .
\end{equation}
Hence \(\phi=0\) is a maximum, not a minimum.  Linear perturbations satisfy
\[
 r^2+3Hr-\frac{2V_0}{\phi_0^2}=0 ,
\]
and the larger root is positive for every \(V_0>0\).  Hubble friction can
delay the roll but cannot make the hilltop asymptotically stable.  The earlier
stable-de-Sitter theorem and oscillation forecast are withdrawn.

For a canonical scalar and pressureless matter, define
\[
 x=\frac{\dot\phi}{\sqrt6\,HM_{\rm Pl}},\qquad
 y=\frac{\sqrt{V}}{\sqrt3\,HM_{\rm Pl}},\qquad
 \lambda=-M_{\rm Pl}\frac{V'}{V},\qquad N=\ln a .
\]
The correct autonomous system is
\begin{align}
 x'&=-3x+\sqrt{\frac32}\lambda y^2
      +\frac32x(1+x^2-y^2),\label{eq:dx}\\
 y'&=-\sqrt{\frac32}\lambda xy
      +\frac32y(1+x^2-y^2).\label{eq:dy}
\end{align}
With \(\Omega_\phi=x^2+y^2\) and
\(w_\phi=(x^2-y^2)/(x^2+y^2)\),
\begin{equation}\label{eq:dOmega-exact}
 \frac{\dd\Omega_\phi}{\dd N}
 =-3w_\phi\,\Omega_\phi(1-\Omega_\phi).
\end{equation}
This does not reduce to the tanh equation printed in earlier versions.

A bounded tanh curve may still be used as a data parametrisation if it is
labelled as such.  A dimensionally consistent non-negative choice is
\begin{align}
 X_{\rm DE}^{\rm proxy}(a)
 &=\frac{\Phi_0}{2}\left[1+
 \tanh\!\bigl(\kappa(a-a_t)\bigr)\right],\label{eq:tanh-solution}\\
 \frac{\dd X_{\rm DE}^{\rm proxy}}{\dd a}
 &=2\kappa X_{\rm DE}^{\rm proxy}
 \left(1-\frac{X_{\rm DE}^{\rm proxy}}{\Phi_0}\right).
 \label{eq:saturation-ODE}
\end{align}
Here \(X_{\rm DE}=\rho_{\rm DE}/\rho_{\rm crit,0}\), not the instantaneous
fraction \(\rho_{\rm DE}/\rho_{\rm crit}(a)\).  Equations
\eqref{eq:tanh-solution}--\eqref{eq:saturation-ODE} are a phenomenological
proxy and are not derived from Eq.~\eqref{eq:Lagrangian}.

\section{Kinematic consequences of a density proxy}\label{sec:observations}

The conditional estimate from Eq.~\eqref{eq:bound} is
\(3.8\times10^{-55}\,{\rm m}^{-2}\), not
\(3.8\times10^{-53}\,{\rm m}^{-2}\).  It is therefore about \(290\) times
below \(\Lambda_{\rm obs}\), not a factor of three.

A bounded curve may be used as a phenomenological density history only if its
normalisation is defined consistently.  Let
\(X_{\rm DE}(a)=\rho_{\rm DE}(a)/\rho_{\rm crit,0}\) and
\[
 E^2(a)=\frac{H^2}{H_0^2}=\Omega_ma^{-3}+X_{\rm DE}(a).
\]
Then
\begin{align}
 q(a)&=-1-\frac{a}{2E^2}\frac{\dd E^2}{\dd a}
 =\frac{\tfrac12\Omega_ma^{-3}-X_{\rm DE}
 -\tfrac a2X_{\rm DE}'}{E^2},\label{eq:q-correct}\\
 w_{\rm DE}(a)&=-1-\frac13\frac{\dd\ln X_{\rm DE}}{\dd\ln a}.
 \label{eq:weff}
\end{align}
The earlier deceleration formula used the wrong coefficient for
\(aX_{\rm DE}'\).  It also alternated between an instantaneous density
fraction and a density normalised to today's critical density.  Consequently,
the reported acceleration redshift, \(w_0\), CPL map, and proxy
\(\chi^2\) are withdrawn pending recomputation from a physically derived
history.  No official DESI likelihood evaluation is claimed.

\begin{remark}[DESI DR2 phantom crossing and scalar-tensor reconstructions]
Recent analyses of DESI DR2 data in combination with Planck 2018 and Supernovae Ia suggest an effective dark energy equation of state $w(z)$ that crosses the phantom divide $w=-1$. While minimally coupled single scalar fields (quintessence) cannot cross $w=-1$ without ghost degrees of freedom, reconstructed scalar-tensor gravity theories~\cite{Efstratiou2025} and interacting dark sector models with field-dependent dark matter masses~\cite{Antusch2026,Sahoo2026} naturally generate an apparent phantom divide $w_{\rm eff}(z) < -1$. The phenomenological proxy in Eq.~\eqref{eq:tanh-solution} serves as a kinematic testbed for such effective scalar-tensor evolution.
\end{remark}

\section{Local consistency gates}

\subsection{Radiation-era trace}\label{sec:BBN}

For an ideal perfect fluid with the sign convention used here,
\[
 T=(1-3w)\rho .
\]
Thus \(T=0\) for exactly massless radiation with \(w=1/3\).  This removes the
leading classical trace source, but it is not a proof of BBN safety: massive
species, trace anomalies, background evolution, and initial conditions must
be propagated through a quantitative abundance calculation.

\subsection{Convexity and screening}\label{sec:screening-gamma}

A fixed strictly convex function of \(\rho\) has one minimiser.  That elementary
fact does not constitute a no-go theorem for chameleon screening.  In a
chameleon model the effective scalar potential depends on ambient density, so
the function being minimised changes with the environment.  The previous
\(\Gamma\)-convergence and energy-dissipation block did not supply a valid
liminf/recovery proof and mixed gradients of different variables; it is
withdrawn.  The accompanying numerical phase plot remains, at most, a toy
illustration and not a Cassini calculation.

\subsection{Scope of the variational statement}\label{sec:resolvent-vacuum}

The variational theorem below concerns only the reduced scalar function with
specified positive coefficients.  It does not solve the radiative-stability
problem, determine a vacuum subtraction scheme, construct an RG trajectory,
or derive a local screening profile.  Those are independent gates.  In
particular, neither a resolvent analogy nor a metric on configuration space
changes the zero set of the differential of a fixed scalar functional.

\section{Renormalisation, beta response, and observational status}
\label{sec:renorm}

The physical scaling problem is concentrated in the matching coefficients
\(B\) and \(C\).  A valid derivation must provide a source- and
scheme-controlled ledger
\[
 {\rm bare}\longrightarrow{\rm Minkowski\ subtraction}
 \longrightarrow{\rm counterterms}
 \longrightarrow{\rm horizon\ remainder}.
\]
No such derivation is supplied here.  The RG, data-processing, one-loop, BBN,
and asymptotic-safety passages in the earlier skeleton were analogies or
leading-order observations, not proofs.  In particular, the ideal-radiation
identity \(T=(1-3w)\rho=0\) is correct for \(w=1/3\), but it does not by itself
guarantee sub-percent BBN safety once massive species, trace anomalies, and
the background scalar dynamics are included.

The earlier definition
\[
 \beta_\Lambda^{\rm old}
 =\frac{\dd}{\dd\ln a}
 \left.\frac{\partial\Phi}{\partial\rho}\right|_{\rho=\rho_\ast(a)}
\]
is identically zero because the bracket is zero for every \(a\).  It cannot
yield a non-zero flow.  If one merely differentiates the conditional ansatz
\(\rho_\ast^{\rm ansatz}\propto H^2M_{\rm Pl}^2/L\),
\(L=\ln(M_{\rm Pl}/H)\), the result is
\begin{equation}\label{eq:beta-response}
 \frac{\dd\rho_\ast^{\rm ansatz}}{\dd\ln a}
 =\rho_\ast^{\rm ansatz}
 \left(2+\frac1L\right)\frac{\dd\ln H}{\dd\ln a}.
\end{equation}
This is a response of an assumed scaling, not an independently extracted beta
function and not evidence for that scaling.

No Wilsonian beta function, fixed-point monotonicity, bound on \(\gamma\) at one loop,
or Kullback--Leibler data-processing statement is promoted to a result in this
revision. Each requires a dimensionally complete action, normalised
probability objects where information inequalities are invoked, and an explicit
renormalisation prescription.

The conditional scale
\[
 \Leff^{\rm ansatz}\simeq3.8\times10^{-55}\ {\rm m}^{-2}
\]
is about \(290\) times below the observed value.  Accordingly, this paper does
not claim a solution of the cosmological-constant or coincidence problems,
a DESI-compatible equation of state, or a stable late-time attractor.

\subsection{\texorpdfstring{Metric \(f(R)\) perturbations}{Metric f(R) perturbations}}\label{sec:predictions-falsifiable}

Use the convention
\[
 \dd s^2=-(1+2\Psi)\dd t^2+a^2(1-2\Phi)\dd\mathbf x^2
\]
and define
\(s(k)= (k^2/a^2)/(k^2/a^2+m_s^2)\).
For \(F=\dd f/\dd R\simeq1\), the quasistatic linear response has the standard
form
\begin{align}
 k^2\Psi&=-4\pi Ga^2\bar\rho\delta
 \left(1+\frac{s}{3}\right),\label{eq:Poisson-Psi}\\
 k^2\Phi&=-4\pi Ga^2\bar\rho\delta
 \left(1-\frac{s}{3}\right).\label{eq:Poisson-Phi}
\end{align}
Thus, with the standard numerator convention,
\begin{equation}\label{eq:grav-slip}
 \eta_{\rm std}=\frac{\Phi}{\Psi}
 =\frac{1-s/3}{1+s/3}
 \longrightarrow\frac12
 \quad (s\to1),
\end{equation}
while the reciprocal convention tends to \(2\).  The Weyl/lensing
combination remains GR-like, apart from the usual \(1/F\) factor:
\begin{equation}\label{eq:Sigma}
 \Sigma\simeq\frac1F\simeq1 .
\end{equation}
Consequently, \(\eta=5/4\) and \(\Sigma=9/8\) are not predictions of standard
metric \(f(R)\) gravity and all forecasts based on those values are withdrawn.
The \(4/3\) enhancement applies to the force potential \(\Psi\), not to a
universal free-energy coefficient~\cite{HuSawicki2007,SotiriouFaraoni2010}.

Strict convexity of a fixed, environment-independent scalar function
\(\Phi_{\rm red}(\rho)\) excludes a second minimum of that same function.
It does not rule out chameleon screening, because the effective scalar
potential itself changes with ambient density.  Likewise, changing only a
positive-definite metric on configuration space cannot move the minimiser of
a fixed scalar function: \({\rm grad}_g\Phi=0\) if and only if
\(\dd\Phi=0\).

\subsection{Source-bound local gate}\label{sec:mcmc-results}

The 2 July 2026 environment ledger uses
\begin{equation}\label{eq:environment-contract}
 |f_R({\rm env})|
 =|f_{R0}|\left(\frac{R_0}{R_{\rm env}}\right)^{n+1}
\end{equation}
together with
\[
 \epsilon_{\rm th}
 =\frac{|f_{R,{\rm env}}-f_{R,{\rm in}}|}{2\Phi_\odot},
 \qquad f_{R,{\rm in}}\simeq0,\quad
 \Phi_\odot=2.12\times10^{-6}.
\]
The Cassini threshold \(\epsilon_{\rm th}<2.3\times10^{-5}\) requires
\begin{equation}
 \frac{|f_R({\rm env})|}{|f_{R0}|}
 <\frac{2\Phi_\odot(2.3\times10^{-5})}{|f_{R0}|}.
\end{equation}
For \(n=1\), this gives
\[
 \frac{R_{\rm env}}{R_0}\ge350.787
 \quad\text{for }|f_{R0}|=1.2\times10^{-5},
 \qquad
 \frac{R_{\rm env}}{R_0}\ge101.264
 \quad\text{for }|f_{R0}|=10^{-6}.
\]
These threshold calculations are arithmetically correct.  No independent
local density/curvature model currently supplies the required
\(R_{\rm env}/R_0\); the ledger therefore permits no claim upgrade.
Neither the exploratory best-fit \(1.2\times10^{-5}\) nor the conservative
\(10^{-6}\) guide is by itself a Cassini proof.  The former is
background-proxy compatible only and is not demonstrated Cassini-safe.

The previous exploratory MCMC scan used simplified diagonal summary terms
rather than the official Planck/DESI likelihoods and encoded Cassini through a
linear hard proxy.  Its numerical samples may be retained as a background
stress test only; they do not establish model preference, official parameter
constraints, or local screening.

\section{Proof status and next steps}\label{sec:discussion}

\begin{center}
\begin{tabularx}{\linewidth}{@{}lY@{}}
\toprule
\textbf{Item} & \textbf{Status after the 16 July 2026 audit}\\
\midrule
Reduced functional & Proven: existence, uniqueness, strict convexity, and the Lambert-\(W\) minimiser for specified \(B,C>0\).\\
Raw mode sum & Invalid as a dimensionally homogeneous functional; motivation only.\\
Physical \(H_0^2/\ln\) scaling & Open and conditional on explicit RG matching; current arithmetic misses observation by about \(290\).\\
P\"oschl--Teller dynamics & Hilltop transient; no stable de-Sitter attractor and no derived tanh law.\\
Metric \(f(R)\) slip/lensing & Standard unscreened limits are \(\eta_{\rm std}\to1/2\) (or reciprocal \(2\)) and \(\Sigma\simeq1/F\), not \(5/4\) and \(9/8\).\\
Local screening & Open: threshold ledger correct, but all numerical passes are source-blocked.\\
Exploratory MCMC & Background-level proxy only; no official Planck/DESI likelihood and no Cassini-safe classification.\\
\bottomrule
\end{tabularx}
\end{center}

The next decisive calculations are: (i) determine dimensionally explicit,
renormalised \(B(H)\) and \(C(H)\); (ii) replace the hilltop potential or
derive a stable scalar history from the full equations; (iii) supply an
independent galactic/solar Hu--Sawicki profile whose curvature ratio meets the
ledger threshold; and (iv) only then rerun an official cosmological
likelihood.  No release or claim upgrade is warranted before those gates are
closed.

\clearpage

\subsection*{AI Disclosure}
In the preparation of this manuscript, AI-assisted tools (Claude, Anthropic)
were used in a supporting capacity, particularly for literature research,
text structuring, and formulation assistance.  The conceptual design,
scientific argumentation, and responsibility for the entire content
rest solely with the author.

\begin{thebibliography}{99}

\bibitem{Weinberg1989}
S.~Weinberg,
\textit{The cosmological constant problem},
Rev.\ Mod.\ Phys.\ \textbf{61}, 1--23 (1989).
\href{https://doi.org/10.1103/RevModPhys.61.1}{doi:10.1103/RevModPhys.61.1}.

\bibitem{Zeldovich1968}
Ya.~B.~Zeldovich,
\textit{The cosmological constant and the theory of elementary particles},
Sov.\ Phys.\ Usp.\ \textbf{11}, 381--393 (1968).
\href{https://doi.org/10.1070/PU1968v011n03ABEH003927}{doi:10.1070/PU1968v011n03ABEH003927}.

\bibitem{PeeblesRatra2003}
P.~J.~E.~Peebles and B.~Ratra,
\textit{The cosmological constant and dark energy},
Rev.\ Mod.\ Phys.\ \textbf{75}, 559--606 (2003).
\href{https://doi.org/10.1103/RevModPhys.75.559}{doi:10.1103/RevModPhys.75.559}.

\bibitem{Planck2018}
Planck Collaboration (N.~Aghanim et al.),
\textit{Planck 2018 results. VI. Cosmological parameters},
Astron.\ Astrophys.\ \textbf{641}, A6 (2020).
\href{https://doi.org/10.1051/0004-6361/201833910}{doi:10.1051/0004-6361/201833910}.

\bibitem{Carroll2001}
S.~M.~Carroll,
\textit{The cosmological constant},
Living Rev.\ Relativ.\ \textbf{4}, 1 (2001).
\href{https://doi.org/10.12942/lrr-2001-1}{doi:10.12942/lrr-2001-1}.

\bibitem{Riess1998}
A.~G.~Riess et al.,
\textit{Observational evidence from supernovae for an accelerating
universe and a cosmological constant},
Astron.\ J.\ \textbf{116}, 1009--1038 (1998).
\href{https://doi.org/10.1086/300499}{doi:10.1086/300499}.

\bibitem{Perlmutter1999}
S.~Perlmutter et al.,
\textit{Measurements of $\Omega$ and $\Lambda$ from 42 high-redshift
supernovae},
Astrophys.\ J.\ \textbf{517}, 565--586 (1999).
\href{https://doi.org/10.1086/307221}{doi:10.1086/307221}.

\bibitem{MartyushevSeleznev2006}
L.~M.~Martyushev and V.~D.~Seleznev,
\textit{Maximum entropy production principle in physics, chemistry and biology},
Phys.\ Rep.\ \textbf{426}, 1--45 (2006).
\href{https://doi.org/10.1016/j.physrep.2005.12.001}{doi:10.1016/j.physrep.2005.12.001}.

\bibitem{Starobinsky1980}
A.~A.~Starobinsky,
\textit{A new type of isotropic cosmological models without singularity},
Phys.\ Lett.\ B \textbf{91}, 99--102 (1980).
\href{https://doi.org/10.1016/0370-2693(80)90670-X}{doi:10.1016/0370-2693(80)90670-X}.

\bibitem{DESI2025}
DESI Collaboration (M.~Abdul-Karim et al.),
\textit{DESI DR2 Results II: Measurements of Baryon Acoustic Oscillations and Cosmological Constraints},
Phys.\ Rev.\ D \textbf{112}, 083515 (2025).
\href{https://doi.org/10.1103/tr6y-kpc6}{doi:10.1103/tr6y-kpc6}.
arXiv:2503.14738 [astro-ph.CO].

\bibitem{Eichhorn2020}
A.~Eichhorn,
\textit{Asymptotically safe gravity},
arXiv:2003.00044 (2020).

\bibitem{EichhornHeld2018}
A.~Eichhorn and A.~Held,
\textit{Top mass from asymptotic safety},
Phys.\ Lett.\ B \textbf{777}, 217--221 (2018).
\href{https://doi.org/10.1016/j.physletb.2017.12.040}{doi:10.1016/j.physletb.2017.12.040}.
arXiv:1707.01107.

\bibitem{HuSawicki2007}
W.~Hu and I.~Sawicki,
\textit{Models of $f(R)$ cosmic acceleration that evade solar-system tests},
Phys.\ Rev.\ D \textbf{76}, 064004 (2007).
\href{https://doi.org/10.1103/PhysRevD.76.064004}{doi:10.1103/PhysRevD.76.064004}.

\bibitem{SotiriouFaraoni2010}
T.~P.~Sotiriou and V.~Faraoni,
\textit{$f(R)$ theories of gravity},
Rev.\ Mod.\ Phys.\ \textbf{82}, 451--497 (2010).
\href{https://doi.org/10.1103/RevModPhys.82.451}{doi:10.1103/RevModPhys.82.451}.

\bibitem{BirrellDavies1982}
N.~D.~Birrell and P.~C.~W.~Davies,
\textit{Quantum Fields in Curved Space},
Cambridge Monographs on Mathematical Physics, Cambridge University Press, 1982.
\href{https://doi.org/10.1017/CBO9780511622632}{doi:10.1017/CBO9780511622632}.

\bibitem{MONDEntropy2025}
{\"O}.~Sevin{\c c}, {\"O}.~{\"O}kc{\"u}, and E.~Aydiner,
\textit{From MOND entropy to extended uncertainty principles: A unified framework},
Phys.\ Lett.\ B \textbf{873}, 140169 (2026).
\href{https://doi.org/10.1016/j.physletb.2026.140169}{doi:10.1016/j.physletb.2026.140169}.
arXiv:2601.14353 [gr-qc].

\bibitem{FinslerJCAP2025}
C.~Pfeifer, N.~Voicu, A.~Friedl-Sz\'asz, and E.~Popovici-Popescu,
\textit{From kinetic gases to an exponentially expanding universe -- the Finsler--Friedmann equation},
JCAP \textbf{2025}(10), 050 (2025).
\href{https://doi.org/10.1088/1475-7516/2025/10/050}{doi:10.1088/1475-7516/2025/10/050}.

\bibitem{HuntMOND2026}
T.~A.~Hunt,
\textit{Dark Matter Phenomenology as Informational Persistence: Nonlocal Gravity, MOND, and the Survival of Curvature Under Coarse-Graining},
Zenodo preprint (2026).
\href{https://doi.org/10.5281/zenodo.18215660}{doi:10.5281/zenodo.18215660}.

\bibitem{GeigerRH2026c}
L.~Geiger, \emph{From Landscape to Proof: The Even-Dominance Route to the Riemann Hypothesis}, Zenodo preprint, 2026.
\href{https://doi.org/10.5281/zenodo.19035640}{doi:10.5281/zenodo.19035640}.

\bibitem{FST-Unified2026}
L.~Geiger, \emph{FST Spectrum Duality: The Renormalized Free Energy Principle as Physical Instantiation of Functional Stability Theory}, Zenodo preprint, 2026.
\href{https://doi.org/10.5281/zenodo.19036190}{doi:10.5281/zenodo.19036190}.

\bibitem{FST-NS2026}
L.~Geiger, \emph{A Thermodynamic Obstruction to Finite-Time Blow-Up in the 3D Navier--Stokes Equations},
Zenodo preprint, 2026.
\href{https://doi.org/10.5281/zenodo.19087449}{doi:10.5281/zenodo.19087449}.

\bibitem{Geiger2026-YM}
L.~Geiger, \emph{Volume-Independent Spectral Gap for Strong-Coupling Lattice Gauge Theory via Poincar\'e--Dobrushin Machinery},
Zenodo preprint, 2026.
\href{https://doi.org/10.5281/zenodo.19087433}{doi:10.5281/zenodo.19087433}.

\bibitem{Geiger2026-BSD}
L.~Geiger, \emph{The BSD Conjecture as a Positivity Normal-Form Theorem},
Zenodo preprint, 2026.
\href{https://doi.org/10.5281/zenodo.19087443}{doi:10.5281/zenodo.19087443}.

\bibitem{CopelandLiddleWands1998}
E.~J.~Copeland, A.~R.~Liddle, and D.~Wands,
\textit{Exponential potentials and cosmological scaling solutions},
Phys.\ Rev.\ D \textbf{57}, 4686--4690 (1998).
\href{https://doi.org/10.1103/PhysRevD.57.4686}{doi:10.1103/PhysRevD.57.4686}.

\bibitem{Li2004}
M.~Li,
\textit{A model of holographic dark energy},
Phys.\ Lett.\ B \textbf{603}, 1--5 (2004).
\href{https://doi.org/10.1016/j.physletb.2004.10.014}{doi:10.1016/j.physletb.2004.10.014}.

\bibitem{BoussoPolchinski2000}
R.~Bousso and J.~Polchinski,
\textit{Quantization of four-form fluxes and dynamical neutralization of the cosmological constant},
JHEP \textbf{0006}, 006 (2000).
\href{https://doi.org/10.1088/1126-6708/2000/06/006}{doi:10.1088/1126-6708/2000/06/006}.

\bibitem{Susskind2003}
L.~Susskind,
\textit{The anthropic landscape of string theory},
in: \textit{Universe or Multiverse?}, Cambridge University Press, 247--266 (2007).
\href{https://doi.org/10.1017/CBO9781107050990.018}{doi:10.1017/CBO9781107050990.018}.

\bibitem{KaloperPadilla2014}
N.~Kaloper and A.~Padilla,
\textit{Sequestering the standard model vacuum energy},
Phys.\ Rev.\ Lett.\ \textbf{112}, 091304 (2014).
\href{https://doi.org/10.1103/PhysRevLett.112.091304}{doi:10.1103/PhysRevLett.112.091304}.

\bibitem{LinYang2004}
F.~Lin and Y.~Yang,
\textit{Existence of energy minimizers as stable knotted solitons in the Faddeev model},
Commun.\ Math.\ Phys.\ \textbf{249}, 273--303 (2004).
\href{https://doi.org/10.1007/s00220-004-1110-y}{doi:10.1007/s00220-004-1110-y}.

\bibitem{Wetterich2024}
C.~Wetterich,
\textit{Dark energy evolution from quantum gravity},
arXiv:2407.03465 [hep-th] (2024).

\bibitem{VasakCCGG2020}
D.~Vasak, J.~Struckmeier, and J.~Kirsch,
\textit{Dark energy and inflation invoked in CCGG by locally contorted space-time},
Eur.\ Phys.\ J.\ Plus \textbf{135}, 404 (2020).
\href{https://doi.org/10.1140/epjp/s13360-020-00415-7}{doi:10.1140/epjp/s13360-020-00415-7}.

\bibitem{ShapiroSola2009}
I.~L.~Shapiro and J.~Sol\`{a},
\textit{The scaling evolution of the cosmological constant},
JHEP \textbf{2002}(02), 006 (2002).
\href{https://doi.org/10.1088/1126-6708/2002/02/006}{doi:10.1088/1126-6708/2002/02/006}.
arXiv:hep-th/0012227.

\bibitem{SolaPeracaula2022}
J.~Sol\`{a} Peracaula,
\textit{The cosmological constant problem and running vacuum in the expanding universe},
Phil.\ Trans.\ R.\ Soc.\ A \textbf{380}, 20210182 (2022).
\href{https://doi.org/10.1098/rsta.2021.0182}{doi:10.1098/rsta.2021.0182}.

\bibitem{CohenKaplanNelson1999}
A.~G.~Cohen, D.~B.~Kaplan, and A.~E.~Nelson,
\textit{Effective Field Theory, Black Holes, and the Cosmological Constant},
Phys.\ Rev.\ Lett.\ \textbf{82}, 4971 (1999).
\href{https://doi.org/10.1103/PhysRevLett.82.4971}{doi:10.1103/PhysRevLett.82.4971}.

\bibitem{GeigerFST-Hub-2026}
L.~Geiger,
\textit{Functional Stability Theory --- A Programmatic Hub: Universal Convexity-Uniqueness Across Scales (From Particles to Cosmology)},
Zenodo (2026).
\href{https://doi.org/10.5281/zenodo.20130499}{doi:10.5281/zenodo.20130499}.

\bibitem{Efstratiou2025}
D.~Efstratiou, E.~A.~Paraskevas, and L.~Perivolaropoulos,
\textit{Addressing the DESI DR2 Phantom-Crossing Anomaly and Enhanced $H_0$ Tension with Reconstructed Scalar-Tensor Gravity},
arXiv:2511.04610 [astro-ph.CO] (2025).

\bibitem{Antusch2026}
S.~Antusch, S.~F.~King, and X.~Wang,
\textit{Coupled Dark Energy and Dark Matter for DESI: An Effective Guide to the Phantom Divide},
arXiv:2604.08449 [hep-ph] (2026).

\bibitem{Sahoo2026}
P.~Thanankullaphong, P.~Sahoo, P.~H.~Puttasiddappa, and N.~Roy,
\textit{Quintom Dark Energy: Future Attractor and Phantom Crossing in Light of DESI DR2 Observation},
arXiv:2601.02284 [gr-qc] (2026).

\end{thebibliography}

\end{document}
